% =============================================================================
% Cahier de calculs manuels du SCR — support pédagogique
% Compilation : pdflatex cahier_calculs.tex (deux fois)
% Les chiffres proviennent de chiffres_pedago.tex, engendré par generer_chiffres.py
% =============================================================================
\documentclass[11pt,a4paper]{article}

\usepackage[T1]{fontenc}
\usepackage{textcomp}
\usepackage[utf8]{inputenc}
\IfFileExists{french.ldf}{\usepackage[french]{babel}}{%
  \renewcommand{\tablename}{Tableau}\renewcommand{\figurename}{Figure}
  \renewcommand{\contentsname}{Table des matières}}
\IfFileExists{lmodern.sty}{\usepackage{lmodern}}{\usepackage{mathptmx}}
\usepackage[margin=2.3cm,headheight=14pt]{geometry}
\usepackage{amsmath}
\usepackage{booktabs}
\usepackage{tabularx}
\usepackage{xcolor}
\usepackage{siunitx}
\usepackage{fancyhdr}
\usepackage{titlesec}
\usepackage{enumitem}
\usepackage{microtype}
\usepackage[hidelinks]{hyperref}
\usepackage{mdframed}

\sisetup{group-separator={\,},group-minimum-digits=4,output-decimal-marker={,}}
\definecolor{bleufonce}{HTML}{1F3864}
\definecolor{grisclair}{HTML}{F4F6FA}
\definecolor{vertfonce}{HTML}{375623}
\definecolor{rouge}{HTML}{C00000}

\titleformat{\section}{\color{bleufonce}\large\bfseries}{\thesection}{0.7em}{}
\titleformat{\subsection}{\color{bleufonce}\normalsize\bfseries}{\thesubsection}{0.6em}{}
\setlist[itemize]{itemsep=2pt,topsep=3pt,leftmargin=1.2em}
\setlist[enumerate]{itemsep=2pt,topsep=3pt,leftmargin=1.5em}

\pagestyle{fancy}\fancyhf{}
\fancyhead[L]{\small\color{bleufonce}Cahier de calculs manuels du SCR}
\fancyhead[R]{\small\color{bleufonce}\thepage}
\renewcommand{\headrulewidth}{0.4pt}

\newmdenv[backgroundcolor=grisclair,linecolor=bleufonce,linewidth=0.8pt,
          innertopmargin=6pt,innerbottommargin=6pt,skipabove=8pt,skipbelow=8pt]{methode}
\newcommand{\piege}[1]{\par\smallskip\noindent\textcolor{rouge}{\textbf{Piège.}} #1\par\smallskip}
\newcommand{\astuce}[1]{\par\smallskip\noindent\textcolor{vertfonce}{\textbf{À retenir.}} #1\par\smallskip}
\newcommand{\Meuro}{\text{M\texteuro}}

\input{chiffres_pedago}

\begin{document}

\begin{titlepage}
\centering
\vspace*{3cm}
{\color{bleufonce}\Huge\bfseries Cahier de calculs manuels du SCR\par}
\vspace{0.8cm}
{\LARGE Comprendre la formule standard en la refaisant à la main\par}
\vspace{2cm}
\begin{minipage}{0.82\linewidth}\small
Ce cahier reprend, module par module, les calculs du projet « compagnie composite ». Chaque
notion est d'abord posée en formule, puis appliquée à un chiffre réel du portefeuille, avec
toutes les étapes intermédiaires. L'objectif est qu'un lecteur qui découvre Solvabilité~II
puisse refaire chaque calcul avec une calculatrice, et retrouver exactement les résultats du
modèle.

\vspace{0.4cm}
Compagnie : actif de \num{11000} \Meuro{}, date d'évaluation 31 décembre 2025, courbe des taux
publiée par l'EIOPA. Tous les montants sont en millions d'euros.
\end{minipage}
\vfill
\end{titlepage}

\tableofcontents
\newpage

% =============================================================================
\section{L'idée centrale, en une page}

Solvabilité~II demande à l'assureur de détenir assez de fonds propres pour absorber un choc
bicentenaire. Formellement, le SCR est la valeur en risque à \qty{99.5}{\percent}, à l'horizon
d'un an, de la variation des fonds propres économiques.

La formule standard remplace cette définition probabiliste par une recette : appliquer des chocs
calibrés, mesurer la perte de fonds propres, puis agréger. D'où l'unique formule à retenir :

\begin{methode}
\[
\boxed{\ \mathrm{SCR}_{\text{sous-module}} \;=\; \max\bigl(0,\; \mathrm{FPB}_{\text{avant}}
- \mathrm{FPB}_{\text{après choc}}\bigr)\ }
\]
où les fonds propres de base valent
\[
\mathrm{FPB} \;=\; \underbrace{A}_{\text{actif en valeur de marché}}
\;-\; \underbrace{\mathrm{BE}}_{\text{best estimate}}
\;-\; \underbrace{\mathrm{MR}}_{\text{marge de risque}}
\;-\; \text{autres passifs}.
\]
\end{methode}

\piege{Le SCR n'est \emph{pas} la perte d'actif. Si l'actif baisse de 100 et que les provisions
baissent de 30 parce que la participation aux bénéfices s'ajuste, le SCR vaut 70. C'est toute la
différence entre calcul brut et calcul net, et c'est la source d'erreur numéro un.}

Trois familles de chocs coexistent :
\begin{itemize}
\item les chocs \textbf{de marché}, qui frappent d'abord l'actif (taux, action, immobilier,
      spread, change, concentration) ;
\item le choc \textbf{de contrepartie}, qui suit une logique de perte en cas de défaut ;
\item les chocs \textbf{de souscription}, qui frappent le passif en modifiant les hypothèses
      (mortalité, rachats, frais, sinistralité).
\end{itemize}

L'agrégation se fait ensuite par des matrices de corrélation, toujours avec la même forme :
\[
\mathrm{SCR}_{\text{agrégé}} = \sqrt{\sum_{i}\sum_{j} \mathrm{Corr}_{i,j}\,
\mathrm{SCR}_i\,\mathrm{SCR}_j}.
\]

\astuce{Cette racine est toujours inférieure à la somme des SCR : c'est le bénéfice de
diversification. Deux risques parfaitement corrélés ($\mathrm{Corr}=1$) s'additionnent ; deux
risques indépendants ($\mathrm{Corr}=0$) s'agrègent en racine de la somme des carrés.}

% =============================================================================
\section{Préalable : actualiser, et valoriser une obligation}

Tout part de la courbe des taux sans risque. Un flux $F$ reçu dans $t$ années vaut aujourd'hui
$F \times (1+s_t)^{-t}$, où $s_t$ est le taux zéro-coupon de maturité $t$.

Pour une obligation, on actualise au taux sans risque \emph{majoré du spread} de l'émetteur :
\[
P = \sum_{t=1}^{n} \frac{C}{(1+s_t+\sigma)^t} + \frac{N}{(1+s_n+\sigma)^n}.
\]

\subsection{Application : la ligne \OblId{}}

Caractéristiques : émetteur \OblEmetteur{}, échelon de qualité de crédit \OblCqs{}, coupon
\OblCoupon{}~\%, maturité \OblMaturite{} ans, spread \OblSpread{} points de base, nominal
\OblNominal{} \Meuro{}.

Au bout de \OblMaturite{} ans, le taux sans risque vaut \OblSpotHuit{}~\%. En ajoutant le spread,
on actualise donc à \OblTauxActu{}~\%.

\begin{table}[h]\centering\small
\caption{Les premiers flux, pour 100 de nominal}
\begin{tabular}{crrrr}
\toprule
$t$ & Flux & $s_t$ & $(1+s_t+\sigma)^{-t}$ & Flux actualisé\\
\midrule
\tableauFlux
\bottomrule
\end{tabular}
\end{table}

En sommant sur les \OblMaturite{} années :
\[
P = \underbrace{\OblPVCoupons{}}_{\text{coupons}} + \underbrace{\OblPVPrincipal{}}_{\text{remboursement}}
= \OblPrix{} \quad \text{pour 100 de nominal.}
\]
La valeur de marché de la ligne vaut donc
$\OblNominal{} \times \OblPrix{} / 100 = \OblVM{}$ \Meuro{}.

\astuce{Le prix dépasse 100 parce que le coupon (\OblCoupon{}~\%) est supérieur au rendement
actuel (\OblTauxActu{}~\%) : le titre a été émis quand les taux étaient plus élevés. L'inverse
donne un titre sous le pair, donc une moins-value latente.}

\subsection{La duration modifiée}

La duration modifiée mesure la sensibilité du prix à une variation de taux :
\[
D_m \simeq -\frac{1}{P}\frac{\Delta P}{\Delta s}.
\]
Pour notre ligne, $D_m = \OblDuration{}$. Concrètement, une hausse de 1~\% des taux fait perdre
environ \OblDuration{}~\% de la valeur. On l'utilise dans le sous-module spread, et comme
approximation de contrôle dans le sous-module taux.

% =============================================================================
\section{Risque de marché}

\subsection{Spread (art. 176)}

Le choc de spread est un pourcentage de perte appliqué à la valeur de marché, fonction de
l'échelon de crédit et de la duration :
\[
\mathrm{stress}(q,d) = a_{q,k} + b_{q,k}\,\bigl(d - \underline{d}_k\bigr),
\]
où $k$ est la bande de duration : $d \le 5$, puis $]5,10]$, $]10,15]$, $]15,20]$, $>20$.

\begin{methode}
\textbf{Application.} Notre ligne est en échelon \OblCqs{} avec une duration de \OblDuration{},
donc dans la bande $]5,10]$, où $a = \OblBandeA{}~\%$ et $b = \OblBandeB{}~\%$ :
\[
\mathrm{stress} = \OblBandeA{}\,\% + \OblBandeB{}\,\% \times (\OblDuration{} - 5) = \OblStress{}\,\%
\]
\[
\text{perte} = \OblStress{}\,\% \times \OblVM{} = \OblPerteSpread{}~\Meuro{}
\]
\end{methode}

Le calcul est répété ligne à ligne, puis les pertes sont sommées. Deux exemptions à connaître :
les obligations souveraines de l'Espace économique européen libellées en monnaie locale ont un
facteur nul, et les obligations sécurisées bénéficient d'un traitement réduit.

\piege{La duration est \emph{plancher\-née à 1 an}. Un titre à trois mois en échelon 3 subit donc
le stress d'un titre à un an, pas le quart.}

\subsection{Taux d'intérêt (art. 166 et 167)}

Deux scénarios sont testés, et le plus défavorable est retenu :
\[
s_t^{\uparrow} = s_t + \max\bigl(1\,\%,\; f^{\uparrow}_t \times |s_t|\bigr),
\qquad
s_t^{\downarrow} = \begin{cases} s_t & \text{si } s_t < 0,\\[2pt]
s_t - f^{\downarrow}_t \times |s_t| & \text{sinon.}\end{cases}
\]
Les facteurs $f_t$ décroissent avec la maturité : \num{0.70} à un an, \num{0.42} à dix ans,
\num{0.20} au-delà de 90 ans à la hausse. Ils figurent dans l'onglet \texttt{Shocks} du fichier
publié chaque mois par l'EIOPA.

\begin{methode}
\textbf{Application à notre obligation.} À \OblMaturite{} ans, le taux passe de \OblSpotHuit{}~\%
à \OblSpotHuitChoc{}~\% (soit $+\OblFacteurChoc{}\,\%$ en relatif). En réactualisant tous les flux
avec la courbe choquée :
\[
P^{\uparrow} = \OblPrixChoc{} \quad\Rightarrow\quad
\text{VM}^{\uparrow} = \OblVMChoc{}~\Meuro{}, \qquad
\text{perte} = \OblPerteTaux{}~\Meuro{}.
\]
Contrôle par la duration : $\OblVM{} \times \OblDuration{} \times \Delta s \simeq
\OblPerteApprochee{}$ \Meuro{}. L'écart avec le calcul exact vient de la convexité et du fait que
le choc n'est pas parallèle.
\end{methode}

\piege{Le sous-module taux est le seul où le \emph{passif} bouge autant que l'actif. Une hausse
des taux réduit la valeur des obligations, mais réduit aussi le best estimate puisque les flux
futurs sont actualisés plus fortement. Le SCR est l'effet net sur les fonds propres, souvent bien
plus faible que la perte d'actif.}

\subsection{Actions (art. 168 à 172)}

Trois catégories, trois chocs. L'ajustement symétrique, publié chaque mois, ajoute ou retranche
jusqu'à 10 points selon la position du marché dans son cycle.

\begin{table}[h]\centering\small
\caption{Assiettes et chocs action du portefeuille}
\begin{tabular}{lrrr}
\toprule
Catégorie & Valeur de marché & Choc & Perte\\
\midrule
\tableauAction
\bottomrule
\end{tabular}
\end{table}

L'agrégation entre les deux types se fait avec une corrélation de \ActionRho{} :
\[
\mathrm{SCR}_{\text{action}} = \sqrt{\mathrm{SCR}_1^2 + 2 \times \ActionRho{} \times
\mathrm{SCR}_1 \mathrm{SCR}_2 + \mathrm{SCR}_2^2}
= \sqrt{\ActionUn{}^2 + 2 \times \ActionRho{} \times \ActionUn{} \times \ActionDeux{}
+ \ActionDeux{}^2} = \ActionAgrege{}~\Meuro{}.
\]

\subsection{Immobilier (art. 174)}

Le plus simple des sous-modules : un choc de \ImmoChoc{}~\% sur toute la poche immobilière,
immeubles d'exploitation compris.
\[
\mathrm{SCR}_{\text{immo}} = \ImmoChoc{}\,\% \times \ImmoAssiette{} = \ImmoPerte{}~\Meuro{}.
\]

\subsection{Change (art. 188)}

Pour chaque devise, on teste une hausse et une baisse de \ChangeChoc{}~\% et on retient la pire.

\begin{table}[h]\centering\small
\caption{Positions en devises et perte associée}
\begin{tabular}{lrr}
\toprule
Devise & Exposition & Perte à \ChangeChoc{}~\%\\
\midrule
\tableauChange
\bottomrule
\end{tabular}
\end{table}

Total des expositions hors euro : \ChangeTotal{} \Meuro{}, soit une perte brute de
\ChangePerte{} \Meuro{}. Le SCR retenu par le modèle est inférieur, car une partie de ces actifs
représente des contrats en unités de compte : leur baisse réduit aussi le passif.

\subsection{Concentration (art. 184)}

On ne sanctionne pas l'exposition, mais son \emph{excès} au-delà d'un seuil de tolérance :
\[
\mathrm{XS}_i = \max\Bigl(0,\; \frac{E_i}{A} - \mathrm{CT}_q\Bigr),
\qquad
\mathrm{Conc}_i = A \times \mathrm{XS}_i \times g_q,
\qquad
\mathrm{SCR}_{\text{conc}} = \sqrt{\sum_i \mathrm{Conc}_i^2}.
\]

\begin{methode}
\textbf{Application.} Actif retenu : $A = \ConcActifs{}$ \Meuro{} (hors unités de compte et hors
souverains). Le plus gros groupe d'émetteurs est \ConcEmetteur{}, avec \ConcExpo{} \Meuro{},
soit \ConcPart{}~\% de l'actif. Son échelon donne un seuil de \ConcSeuil{}~\% et un facteur
$g = \ConcFacteur{}~\%$ :
\[
\mathrm{XS} = \ConcPart{}\,\% - \ConcSeuil{}\,\% = \ConcXS{},
\qquad
\mathrm{Conc} = \ConcActifs{} \times \ConcXS{} \times \ConcFacteur{}\,\% = \ConcPremier{}~\Meuro{}.
\]
\end{methode}

\begin{table}[h]\centering\small
\caption{Quatre plus grandes expositions}
\begin{tabular}{lrrcrrr}
\toprule
Groupe & Exposition & Part & Éch. & Seuil & XS & Contribution\\
\midrule
\tableauConcentration
\bottomrule
\end{tabular}
\end{table}

En sommant les carrés sur tous les groupes puis en prenant la racine :
$\mathrm{SCR}_{\text{conc}} = \ConcTotal{}$ \Meuro{}.

\piege{La racine de la somme des carrés suppose les défauts \emph{indépendants} entre émetteurs.
Un seul gros émetteur au-dessus du seuil peut donc porter presque tout le sous-module.}

\subsection{Agrégation du module marché (art. 164)}

\begin{table}[h]\centering\small
\caption{Matrice de corrélation du module marché (scénario de hausse des taux, $A = 0$)}
\begin{tabular}{lcccccc r}
\toprule
& taux & action & immo & spread & change & conc. & SCR\\
\midrule
\tableauMatriceMarche
\bottomrule
\end{tabular}
\end{table}

\[
\mathrm{SCR}_{\text{marché}} = \sqrt{\mathbf{v}^{\top}\mathbf{C}\,\mathbf{v}}
= \sqrt{\MarcheQuadratique{}} = \MarcheAgrege{}~\Meuro{},
\]
à comparer à la somme brute de \MarcheSomme{} \Meuro{} : la diversification économise
\MarcheDiversification{} \Meuro{}.

\astuce{Le paramètre $A$ vaut 0 si c'est le scénario de hausse des taux qui est retenu, et
\num{0.5} si c'est celui de baisse. Autrement dit, une baisse des taux est réputée aller de pair
avec la chute des actions et l'écartement des spreads, alors qu'une hausse est traitée comme
indépendante.}

% =============================================================================
\section{Risque de défaut de la contrepartie}

Ici, la logique change : on ne choque plus un prix, on modélise une perte en cas de défaut.
Deux ingrédients : la perte en cas de défaut (LGD) et la probabilité de défaut (PD), déduite de
l'échelon de qualité de crédit.

\begin{table}[h]\centering\small
\caption{Expositions de type 1}
\begin{tabular}{lcrrr}
\toprule
Contrepartie & Éch. & Exposition & LGD & PD\\
\midrule
\tableauContrepartie
\bottomrule
\end{tabular}
\end{table}

La LGD vaut \qty{50}{\percent} de l'exposition pour les réassureurs — on suppose récupérer la
moitié — et \qty{100}{\percent} pour la trésorerie.

\subsection{La formule de variance}

\[
V = \sum_{j,k} \frac{\mathrm{PD}_j \mathrm{PD}_k}
{1{,}25(\mathrm{PD}_j+\mathrm{PD}_k) - \mathrm{PD}_j\mathrm{PD}_k}\,
\mathrm{TLGD}_j \mathrm{TLGD}_k
\;+\; \sum_j \frac{1{,}5\,\mathrm{PD}_j(1-\mathrm{PD}_j)}{2{,}5-\mathrm{PD}_j}
\sum_{i \in j}\mathrm{LGD}_i^2
\]

Le premier terme capte la dépendance entre contreparties, le second le risque propre à chacune.
On regroupe d'abord les expositions par probabilité de défaut :

\[
\mathrm{TLGD}(\mathrm{PD}=\CpPDun{}\,\%) = \CpTLGDun{}, \qquad
\mathrm{TLGD}(\mathrm{PD}=\CpPDdeux{}\,\%) = \CpTLGDdeux{}.
\]

\begin{methode}
\textbf{Application.} La double somme et le terme diagonal (\CpDiagonale{}) donnent
$V = \CpVariance{}$, donc $\sqrt{V} = \CpRacine{}$ \Meuro{}.

On compare ensuite à la LGD totale, \CpLGDtotal{} \Meuro{} :
\[
\sqrt{V} = \CpRacine{} \;\le\; 5\,\% \times \CpLGDtotal{} = \CpSeuil{}
\quad\Longrightarrow\quad
\mathrm{SCR}_1 = 3\sqrt{V} = \CpUn{}~\Meuro{}.
\]
\end{methode}

\astuce{La règle à trois branches ($3\sqrt{V}$, $5\sqrt{V}$, ou la LGD totale) est une
approximation de quantile : plus le portefeuille est concentré, plus le multiplicateur monte,
jusqu'au cas extrême où l'on provisionne la perte totale.}

\subsection{Type 2 et agrégation}

Les créances suivent un traitement forfaitaire :
\[
\mathrm{SCR}_2 = 15\,\% \times \CpCreances{} + 90\,\% \times \CpEchues{} = \CpDeux{}~\Meuro{},
\]
la forte pénalisation des créances échues depuis plus de trois mois reflétant leur probabilité
de non-recouvrement. Puis
\[
\mathrm{SCR}_{\text{déf}} = \sqrt{\mathrm{SCR}_1^2 + 1{,}5\,\mathrm{SCR}_1\mathrm{SCR}_2
+ \mathrm{SCR}_2^2} = \CpAgrege{}~\Meuro{}.
\]

% =============================================================================
\section{Risque de souscription vie}

Le principe est différent des modules d'actif : on modifie une hypothèse, on \textbf{reprojette
entièrement} les provisions, et le SCR est la hausse du best estimate.

\begin{table}[h]\centering\small
\caption{Les sept chocs du module vie}
\begin{tabularx}{\linewidth}{lXl}
\toprule
Sous-module & Choc & Article\\
\midrule
Mortalité & $q_x \rightarrow 1{,}15\,q_x$, définitivement & 137\\
Longévité & $q_x \rightarrow 0{,}80\,q_x$, définitivement & 138\\
Invalidité & hausse des entrées, baisse des sorties & 139\\
Rachat & $\times 1{,}5$, ou $\times 0{,}5$, ou rachat immédiat de \qty{40}{\percent} & 142\\
Frais & $+10\,\%$ de frais et $+1$ point d'inflation & 143\\
Révision & $+3\,\%$ sur les rentes révisables & 141\\
Catastrophe & $q_x \rightarrow q_x + 0{,}15$ point, la première année seulement & 144\\
\bottomrule
\end{tabularx}
\end{table}

\subsection{Le seul choc calculable de tête : la catastrophe}

La surmortalité de \VieCatChoc{} point s'applique une année, aux capitaux sous risque :
\[
\mathrm{SCR}_{\text{cat}} \simeq \VieCatChoc{}\,\% \times \VieCapitaux{} = \VieCatPerte{}~\Meuro{},
\]
à comparer aux \VieCatModele{} \Meuro{} du modèle. L'écart tient à l'actualisation et au fait que
les décès surviennent en moyenne en milieu d'année.

\astuce{Ce calcul est le meilleur contrôle de cohérence du module vie. Si votre SCR catastrophe
s'écarte de \qty{0.15}{\percent} des capitaux sous risque, c'est que les capitaux sont mal
agrégés.}

\subsection{Mortalité et longévité : pourquoi ils ne s'annulent pas}

Les deux chocs vont en sens inverse, mais ils ne portent pas sur les mêmes contrats. Une hausse
de la mortalité coûte cher sur les garanties décès ; une baisse coûte cher sur les rentes.
Le règlement demande de retenir, pour chaque contrat, le choc qui \emph{augmente} les provisions.
D'où une corrélation négative ($-0{,}25$) dans la matrice, mais jamais une compensation totale.

Dans notre portefeuille, le choc de mortalité coûte \VieMortaliteModele{} \Meuro{}, l'essentiel
venant de la temporaire décès.

\subsection{Rachat massif : le sous-module dominant}

Le choc consiste à supposer que \VieMassifTaux{}~\% des contrats sont rachetés immédiatement.
Pour chaque contrat, on compare :
\begin{itemize}
\item ce que l'on doit payer tout de suite, c'est-à-dire la valeur de rachat ;
\item ce que le contrat valait dans les comptes, c'est-à-dire son best estimate.
\end{itemize}
Si la valeur de rachat dépasse le best estimate, le rachat coûte de l'argent et le choc
s'applique. Sinon, il est favorable et \emph{n'est pas retenu}. Ce tri se fait contrat par
contrat, jamais globalement. Dans notre cas, le rachat ressort à \VieMassifModele{} \Meuro{},
soit le premier poste du module vie.

\piege{Beaucoup d'étudiants appliquent le choc massif à tout le portefeuille et obtiennent un
chiffre trop faible, car les contrats favorables viennent compenser. La sélection contrat par
contrat est explicitement demandée à l'article 142.}

\subsection{Agrégation du module vie}

\begin{table}[h]\centering\small
\caption{Matrice de corrélation vie et SCR par sous-module}
\begin{tabular}{lccccccc r}
\toprule
& mort. & long. & inval. & rachat & frais & révis. & cat. & SCR\\
\midrule
\tableauMatriceVie
\bottomrule
\end{tabular}
\end{table}

Somme brute : \VieSomme{} \Meuro{}. Après agrégation : \VieAgrege{} \Meuro{}.

% =============================================================================
\section{Risque de souscription non-vie}

\subsection{Primes et réserves}

Ce sous-module traite d'un coup deux incertitudes : le coût des sinistres à venir (primes) et
l'écart entre les provisions constituées et le coût final (réserves). Pour chaque ligne
d'activité :
\[
\sigma_s = \frac{\sqrt{(\sigma_{p}V_{p})^2 + \sigma_{p}V_{p}\,\sigma_{r}V_{r} + (\sigma_{r}V_{r})^2}}
{V_{p}+V_{r}}
\]
où $V_p$ est le volume de primes de l'année à venir et $V_r$ le best estimate de sinistres. Les
écarts types $\sigma_p$ et $\sigma_r$ sont donnés par l'annexe~II du règlement.

\begin{methode}
\textbf{Application à la ligne \NvLob{}.} $V_p = \NvVp{}$, $V_r = \NvVr{}$,
$\sigma_p = \NvSigmaP{}~\%$, $\sigma_r = \NvSigmaR{}~\%$ :
\[
\sigma_s = \frac{\NvNumerateur{}}{\NvVp{} + \NvVr{}} = \NvSigmaLob{}\,\%.
\]
\end{methode}

\begin{table}[h]\centering\small
\caption{Volumes et écarts types par ligne d'activité}
\begin{tabular}{lrrrrr}
\toprule
Ligne & $V_p$ & $V_r$ & $\sigma_p$ & $\sigma_r$ & $\sigma_s$\\
\midrule
\tableauNonVie
\bottomrule
\end{tabular}
\end{table}

Les lignes sont ensuite agrégées par la matrice de corrélation de l'annexe~IV, ce qui donne un
écart type global de \NvSigmaGlobal{}~\% sur un volume total de \NvVolumeTotal{} \Meuro{}. Le
SCR vaut alors trois fois cet écart type :
\[
\mathrm{SCR}_{\text{primes, réserves}} = 3 \times \NvSigmaGlobal{}\,\% \times \NvVolumeTotal{}
= \NvPrimesReserves{}~\Meuro{}.
\]

\astuce{Pourquoi 3 ? Parce que sous une hypothèse log-normale, le quantile à \qty{99.5}{\percent}
se situe à environ trois écarts types de la moyenne. Le facteur 3 est donc le
« \num{99.5}\textsuperscript{e} percentile » déguisé.}

\subsection{Catastrophe et réassurance}

La charge catastrophe brute ressort à \NvCatBrut{} \Meuro{}. Le traité en excédent de sinistres,
de priorité \NvPriorite{} \Meuro{} et de portée \NvPortee{} \Meuro{}, la ramène à :
\[
L^{\text{net}} = \min(L, \NvPriorite{}) + \max(0,\; L - \NvPriorite{} - \NvPortee{})
= \NvCatNet{}~\Meuro{}.
\]
Le module non-vie agrégé vaut \NvAgrege{} \Meuro{}.

\piege{La réassurance ne se déduit pas du SCR après coup : elle s'applique sinistre par sinistre,
\emph{avant} l'agrégation. Un traité qui plafonne les sinistres extrêmes réduit surtout la queue
de distribution, donc le SCR, et beaucoup moins la charge moyenne.}

% =============================================================================
\section{Risque opérationnel (art. 204)}

Aucun choc ici : une formule forfaitaire fondée sur les volumes, puis un plafond.
\[
\mathrm{SCR}_{\text{op}} = \min\bigl(30\,\% \times \mathrm{BSCR},\;
\max(\mathrm{Op}_{\text{primes}}, \mathrm{Op}_{\text{provisions}})\bigr)
+ 25\,\% \times \text{frais sur unités de compte}
\]

\begin{methode}
\textbf{Application.}
\begin{align*}
\mathrm{Op}_{\text{primes}} &= 4\,\% \times \OpPrimesVie{} + 3\,\% \times \OpPrimesNonVie{}
= \OpPrimes{}~\Meuro{}\\
\mathrm{Op}_{\text{provisions}} &= 0{,}45\,\% \times \OpProvVie{} + 3\,\% \times \OpProvNonVie{}
= \OpProvisions{}~\Meuro{}\\
\text{maximum retenu} &= \OpProvisions{}, \qquad
\text{plafond} = 30\,\% \times \mathrm{BSCR} = \OpPlafond{}\\
\mathrm{SCR}_{\text{op}} &= \OpProvisions{} + 25\,\% \times \OpFraisUC{} = \OpTotal{}~\Meuro{}
\end{align*}
\end{methode}

\astuce{Le plafond de \qty{30}{\percent} du BSCR ne mord presque jamais pour un assureur vie de
taille normale : il protège surtout contre les portefeuilles à très gros volumes et très faible
risque.}

% =============================================================================
\section{Du BSCR au SCR, puis au MCR}

\subsection{Agrégation des cinq modules}

\begin{table}[h]\centering\small
\caption{Matrice de corrélation du BSCR et SCR par module}
\begin{tabular}{lccccc r}
\toprule
& marché & défaut & vie & santé & non-vie & SCR\\
\midrule
\tableauMatriceBscr
\bottomrule
\end{tabular}
\end{table}

Somme brute : \BscrSomme{} \Meuro{}. Après agrégation : $\mathrm{BSCR} = \BscrAgrege{}$ \Meuro{},
soit \BscrDiversification{} \Meuro{} de diversification.

\subsection{Les deux ajustements}

\textbf{Capacité d'absorption des provisions techniques.} Si, en scénario choqué, l'assureur peut
réduire la participation aux bénéfices futurs, une partie du choc est supportée par les assurés.
L'ajustement est négatif, plafonné par le montant des prestations discrétionnaires futures. Dans
notre modèle déterministe, cette absorption est déjà intégrée dans la reprojection des
provisions : l'ajustement explicite est donc nul. C'est le modèle stochastique qui permet de
l'isoler.

\textbf{Capacité d'absorption des impôts différés.} Une perte bicentenaire réduirait l'impôt
futur. On peut donc déduire cet impôt économisé, à condition de démontrer qu'il serait
effectivement récupérable :
\[
\text{plafond théorique} = 25\,\% \times (\mathrm{BSCR} + \mathrm{SCR}_{\text{op}}) = \LacPlafond{}~\Meuro{}.
\]
La capacité réellement justifiable se limite ici aux \LacIdp{} \Meuro{} d'impôts différés passifs
au bilan, augmentés de \LacPart{}~\% de l'écart, au titre des bénéfices futurs :
$\mathrm{Adj}_{\mathrm{DT}} = \LacDt{}$ \Meuro{}.

\piege{Confondre le plafond réglementaire et la capacité justifiable est l'erreur classique. Le
règlement autorise jusqu'au taux d'impôt appliqué au choc, mais l'entreprise doit prouver qu'elle
retrouvera des bénéfices imposables. Un superviseur regarde d'abord cette démonstration.}

\subsection{Le SCR}

\[
\mathrm{SCR} = \mathrm{BSCR} + \mathrm{SCR}_{\text{op}} + \mathrm{Adj}_{\mathrm{TP}}
+ \mathrm{Adj}_{\mathrm{DT}}
= \BscrAgrege{} + \OpTotal{} + 0 + (\LacDt{}) = \ScrFinal{}~\Meuro{}.
\]

Avec des fonds propres de base de \FondsPropres{} \Meuro{}, le ratio de couverture vaut
\[
\frac{\FondsPropres{}}{\ScrFinal{}} = \RatioScr{}\,\%.
\]

\subsection{Le MCR}

Le minimum de capital requis est un filet de sécurité, calculé de façon linéaire puis encadré :
\begin{align*}
\mathrm{MCR}_{\text{linéaire}} &= \McrVie{} \text{ (vie)} + \McrNonVie{} \text{ (non-vie et santé)}
= \McrLineaire{}~\Meuro{}\\
\text{corridor} &: \max(25\,\% \times \mathrm{SCR}) = \McrPlancher{}, \quad
\min(45\,\% \times \mathrm{SCR}) = \McrPlafond{}\\
\mathrm{MCR} &= \McrRetenu{}~\Meuro{}, \qquad \text{ratio de couverture} = \RatioMcr{}\,\%.
\end{align*}

\astuce{Le MCR se situe toujours entre \qty{25}{\percent} et \qty{45}{\percent} du SCR. S'il sort
de ces bornes, c'est le corridor qui tranche, jamais le calcul linéaire. Franchir le MCR déclenche
le retrait d'agrément : c'est le seuil le plus grave du régime.}

% =============================================================================
\section{Récapitulatif et erreurs fréquentes}

\begin{table}[h]\centering\small
\caption{Le chemin complet, en un tableau}
\begin{tabularx}{\linewidth}{Xr}
\toprule
Étape & \Meuro{}\\
\midrule
Risque de marché (six sous-modules agrégés) & \MarcheAgrege{}\\
Risque de défaut de la contrepartie & \CpAgrege{}\\
Risque de souscription vie & \VieAgrege{}\\
Risque de souscription non-vie & \NvAgrege{}\\
\textbf{BSCR après diversification} & \textbf{\BscrAgrege{}}\\
Risque opérationnel & \OpTotal{}\\
Capacité d'absorption des impôts différés & \LacDt{}\\
\textbf{SCR} & \textbf{\ScrFinal{}}\\
MCR & \McrRetenu{}\\
Fonds propres de base & \FondsPropres{}\\
\textbf{Ratio de couverture du SCR} & \textbf{\RatioScr{}\,\%}\\
\bottomrule
\end{tabularx}
\end{table}

\subsection{Les huit erreurs les plus fréquentes}

\begin{enumerate}
\item Confondre perte d'actif et perte de fonds propres : le passif bouge aussi.
\item Oublier de retenir le maximum entre hausse et baisse des taux, et de choisir le paramètre
      $A$ correspondant.
\item Additionner les sous-modules au lieu de les agréger, ce qui surestime largement le SCR.
\item Appliquer le rachat massif globalement au lieu de contrat par contrat.
\item Retenir un choc favorable : un sous-module ne peut pas produire un SCR négatif.
\item Actualiser les provisions non-vie avec un taux unique au lieu de la courbe et de la cadence
      de règlement.
\item Prendre le plafond réglementaire de la capacité d'absorption fiscale sans justifier des
      bénéfices futurs.
\item Oublier la mise en transparence des OPC : un fonds obligataire porte du risque de spread et
      de taux, pas un risque « fonds ».
\end{enumerate}

\subsection{Pour s'entraîner}

Quelques calculs à refaire seul, dont les réponses figurent plus haut :
\begin{enumerate}
\item Recalculer le prix de l'obligation \OblId{} à partir de la courbe, puis sa perte sous choc
      de spread.
\item Vérifier le SCR catastrophe vie à partir des capitaux sous risque.
\item Refaire l'agrégation du module marché à la main, matrice en main, et retrouver
      \MarcheAgrege{} \Meuro{}.
\item Recalculer $\sigma_s$ pour une autre ligne non-vie que \NvLob{}.
\item Vérifier que le MCR tombe bien dans le corridor \qty{25}{\percent}–\qty{45}{\percent}.
\end{enumerate}

\vfill
\begin{center}\small\itshape
Tous les chiffres de ce cahier sont engendrés depuis les données et les résultats du modèle :
les refaire à la main doit donner exactement les mêmes valeurs.
\end{center}

\end{document}
